% SPDX-License-Identifier: Apache-2.0
\section{Limits}
\label{sec:limitations}

\subsection{One version}

Every claim here was measured on files written by Vivado 2025.2.
Nothing distinguishes a property of the format from a property of that
release, because no second release has been examined. Any use of this
work on another version should begin by re-running the corpus.

\subsection{What is still unexplained}

The repository keeps an open questions section, separate from the
findings table, and it is not short. Among its entries: what the
chunking constants 24 and 299 mean; what several handle space costs
pay for, including the \texttt{0x1088} bytes before the first signal;
which value class codes the corpus has never produced; and what a
numbering added under \texttt{-debug all} is read by. Each is labelled
a guess or a gap. None of them is required by the reader, which is why
they can stay open.

\subsection{The derivation}

An AI agent produced this analysis, under human supervision, by
running experiments and reading bytes. No AMD documentation, source or
binary was used, and the result is not a port, a decompilation or a
specification.

The failure mode of such a derivation is a story that reproduces on
the cases that produced it. The corpus, the truth files, the external
guards and the requirement that every finding name a reproducing
command are all directed at that failure mode. They reduce it; they do
not remove it. Where a wrong answer would be silent and expensive,
the recommendation in the tool's own help text is to open the database
in Vivado instead.

\subsection{Scope of the reader}

The reader decodes the objects a simulation records. It does not
decode the contents of SystemVerilog dynamic objects, which the
database appears to declare and not to record, and it does not write
databases.
